\chapter{Therapeutic Drug Monitoring}
\section*{What Is This Test?} Therapeutic drug monitoring measures selected medicine concentrations to balance efficacy and toxicity.\section*{Alternative Names} Drug-level monitoring; serum drug level; pharmacokinetic monitoring.\section*{Principle} A measured concentration is interpreted with dose, timing, pharmacokinetics, target range, disease, and response.\section*{Why This Test Is Ordered} It is useful for medicines with narrow therapeutic windows, variable clearance, adherence questions, interactions, or toxicity concerns.\section*{Primary vs Secondary Test} Drug concentration is secondary to clinical assessment; a number alone does not determine treatment.\section*{How This Test Is Performed} Blood is collected at a documented time, often immediately before a dose (trough), and measured by immunoassay or mass spectrometry.\section*{What Preparation Is Needed} Record the last dose, dose changes, missed doses, kidney and liver disease, and interacting medicines.\section*{Understanding Results} Timing errors can make a result uninterpretable; target ranges are population guides, not universal goals.\section*{Follow-up Tests} Repeat levels, kidney/liver tests, ECG, toxicity assessment, or pharmacogenetic testing may follow.\section*{References} \begin{enumerate}\item MedlinePlus. Therapeutic Drug Monitoring.\item International Association for Therapeutic Drug Monitoring and Clinical Toxicology. Guidance.\end{enumerate}\clearpage\section*{Understanding Results and Follow-up}
The laboratory report should identify the specimen, method, units, reference interval or decision limit, and whether the result is positive, negative, abnormal, or inconclusive. Reference intervals vary with age, sex, pregnancy, specimen type, assay, and laboratory; a result outside a range is not automatically a diagnosis, and a result within a range does not always exclude disease. Discuss unexpected results with the ordering clinician, who can decide whether repeat or confirmatory testing, treatment, or referral is appropriate. Do not change medicines or delay urgent care based on an isolated result.
\section*{Regulatory Status and Authorization Date}This chapter describes a test category, not a specific commercial product. FDA, CDSCO, EU IVDR, and MHRA approval or registration dates therefore cannot be assigned without the assay manufacturer, product name, instrument, and intended use. For laboratory-developed tests, an individual agency approval date may not exist. See the Preface for the interpretation of regulatory status.
\clearpage
